
\documentclass[aspectratio=169]{beamer}

\mode<presentation>
\usetheme{Boadilla}
\definecolor{NTHU1}{RGB}{127,16,132}
\definecolor{NTHU2}{RGB}{228,210,231}
\definecolor{blue}{RGB}{30,90,205}
\definecolor{red}{RGB}{213,94,0}
\definecolor{green}{RGB}{0,128,0}
\definecolor{charcoalgray}{RGB}{108,104,104}
\setbeamercolor{title}{fg=NTHU1}
\setbeamercolor{frametitle}{fg=NTHU1}
\setbeamercolor{block title}{bg=NTHU1, fg=white}
\setbeamercolor{block body}{bg=NTHU2}
\setbeamercolor{structure}{fg=NTHU1}
\setbeamercolor{item projected}{fg=white}
\setbeamercolor{item}{fg=NTHU1}
\setbeamercolor{subitem}{fg=NTHU1}
\setbeamercolor{section in toc}{fg=NTHU1}
\setbeamercolor{description item}{fg=NTHU1}
\setbeamercolor{caption name}{fg=NTHU1}
\setbeamercolor{button}{bg=NTHU1, fg=white}
\setbeamercolor{caption name}{fg=NTHU1}
\usepackage{graphics}
\usepackage{tikz}
\usepackage{amsmath}
\usepackage{bbm}
\usetikzlibrary{decorations.pathreplacing}
\usepackage{geometry}
\usepackage{booktabs}
\usepackage{bookmark}
\usepackage{multirow, makecell}
\usepackage{float}
\usepackage{fancyvrb}
\usepackage{caption}
\captionsetup[figure]{singlelinecheck = false, format= hang, justification=centering}
\captionsetup[subfigure]{singlelinecheck = false, format= hang, justification=raggedright}
\usepackage{subcaption}
\usepackage{adjustbox}
\usepackage{hyperref}
\usepackage{threeparttable}
\usepackage{appendixnumberbeamer} 
\usepackage[T1]{fontenc}
\usepackage[default]{lato} 
\usepackage{smartdiagram}
\smartdiagramset{
  back arrow disabled = true,
  module minimum width=1cm,
  module x sep = 3,
  uniform arrow color=true,
}
\usepackage{xcolor}
\usepackage{colortbl}
	\definecolor{light-gray}{gray}{0.99}

\newenvironment{wideitemize}{\itemize\addtolength{\itemsep}{10pt}}{\enditemize}
\newenvironment{wideenumerate}{\enumerate\addtolength{\itemsep}{10pt}}{\endenumerate}
\newenvironment{widedescription}{\description\addtolength{\itemsep}{10pt}}{\enddescription}
\hypersetup{
colorlinks=true,
linkcolor=NTHU1,
filecolor=green, 
urlcolor=blue,
}
\beamertemplatenavigationsymbolsempty
\setbeamercolor{author in head/foot}{bg=white, fg=NTHU1}
\setbeamercolor{title in head/foot}{bg=white, fg=NTHU1}
\setbeamercolor{date in head/foot}{bg=white, fg=NTHU1}
\setbeamercolor{section in head/foot}{bg=white, fg=NTHU1}
\setbeamercolor{headline}{bg=white}
\setbeamertemplate{footline}{
    \leavevmode%
    \hbox{%
        \begin{beamercolorbox}[wd=.333333\paperwidth,ht=2.25ex,dp=1ex,center]{date in head/foot}%
            \usebeamerfont{date in head/foot}%\insertdate
        \end{beamercolorbox}%
        \begin{beamercolorbox}[wd=.333333\paperwidth,ht=2.25ex,dp=1ex,center]{title in head/foot}%
            \usebeamerfont{title in head/foot}\insertshorttitle
        \end{beamercolorbox}%
        \begin{beamercolorbox}[wd=.333333\paperwidth,ht=2.25ex,dp=1ex,right]{date in head/foot}%
            \usebeamerfont{date in head/foot}\insertframenumber{}\hspace*{2em}
        \end{beamercolorbox}}%
        \vskip0pt%
    }
    %% May need to script the introduction!!!!!!
%\setbeamercolor{page number in head/foot}{fg=NTHU1}
\setbeamertemplate{section in toc}[sections numbered]
\setbeamertemplate{subsection in toc}{\leavevmode\leftskip=3em\rlap{\hskip-1.75em\inserttocsectionnumber.\inserttocsubsectionnumber}
\inserttocsubsection\par}
\setbeamerfont{subsection in toc}{size=\footnotesize}
%\setbeamertemplate{headline}{%
  %\begin{beamercolorbox}[ht=5.5ex]{section in head/foot}
    %\vskip2pt\insertnavigation{0.33\paperwidth}\vskip2pt
  %\end{beamercolorbox}%
%}
\newenvironment{transitionframe}{\setbeamercolor{background canvas}{bg=white}\setbeamertemplate{footline}{} \begin{frame}[noframenumbering]}{\end{frame}}


\makeatletter
\let\@@magyar@captionfix\relax
\makeatother


\title[]{Public Economics and Development} % Change this regularly
\subtitle[]{Week 14: Monitoring and Incentivizing Public Sector Workers}
\author[Lee]{Seung-hun Lee }
\institute[]{Taipei School of Economics\\ National Tsing Hua University}

\date[]{May 26th, 2026}

\begin{document}
\begin{transitionframe}
\titlepage
\end{transitionframe}


%%% Color slides for section headers: Use for colloquium version (The ones bewteen \iffals and \fi)

\section{Overview setup of the topic}

\begin{frame}
\frametitle{Why politicians face little scrutiny — and why it matters}
 
Citizens hold little information on political activities and limited punishment mechanisms
\begin{itemize}\setlength\itemsep{3pt}
  \item Limited channels (local radio, low-circulation press) to learn about politician behavior
  \begin{itemize}\setlength\itemsep{3pt}
    \item Media and internet are often censored
  \end{itemize}
  \item Actual data on whether and how services reach endpoints is limited
  \item Recall elections and audits can be misused by political rivals and pressure groups
\end{itemize}\medskip
Why accountability still matters
\begin{itemize}\setlength\itemsep{3pt}
  \item Need to prevent diverson of public resources away from intended recipients
  \item Electoral outcomes serve as strong incentives on elected officials
    \begin{itemize}\setlength\itemsep{3pt}
    \item Whereever legally feasible, politicians care about re-election
    \item Even if not, continuation of incumbent party rule still matters
  \end{itemize}
  \item Long run: Shape who enters politics, not just disciplining incumbents
\end{itemize}
 
\end{frame}



\begin{frame}
\frametitle{Media as an accountability tool — promise and limits}
How media promotes accountability
\begin{itemize}\setlength\itemsep{3pt}
  \item Reward well-behaving politicians, not only punishes corrupt ones
  \item Disseminate information via dominant channels (radio, television, web) 
  \item Effective bottom-up enforcement: Campaign for or against various causes
  \item Lower coordination costs and facilitate collective action (esp, social media and mobile)
  
\end{itemize}
 
\medskip
But media can backfire very badly
\begin{itemize}\setlength\itemsep{3pt}
  \item Above positive assumes sufficient media reach (not always the case)
  \item Misinformation can entrench incumbents or manufacture political crises
  \item In low-trust environments, media may not shift voter behavior at all
\end{itemize}\bigskip
$\Rrightarrow$ This is a large field in political economy that deserves its own semester
\end{frame}





\begin{frame}
\frametitle{Incentivizing bureaucrats — types of instruments}
 Financial incentives
\begin{itemize}\setlength\itemsep{3pt}
  \item Nonlinear pay tied to performance (better service delivery and showing up)
  \item Subject to fiscal constraints and politically difficult for entrenched civil servants
 \end{itemize} \medskip
 Nonmonetary measures: Job postings and career progression
 \begin{itemize}\setlength\itemsep{3pt}
  \item Assignment favoring high-performers can induce effort
  \begin{itemize}\setlength\itemsep{3pt}
  \item Potential mismatch: preference of top performance and where they are needed may differ
\end{itemize}
  \item Bureaucrats care about career progression (how fast they are promoted)
   \begin{itemize}\setlength\itemsep{3pt}
  \item Often linked with post-retirement benefits and social prestige
  \item Inflexible trajectory in general: Could lead to them exiting for outside options
  \item So they care about not getting suspensions - may lead to passive behaviors instead
\end{itemize}
\end{itemize}\medskip
Unlike private organizations, role of financial incentives could be limited
\begin{itemize}\setlength\itemsep{3pt}
  \item Payscale is generally inflexible
  \item Side effects at recruitment: Crowd out prosocial individuals (Deserrano 2019)
 \end{itemize}
 
\end{frame}
 






\begin{transitionframe}
\frametitle{Roadmap for today}
\tableofcontents
\end{transitionframe}








\section{Monitoring and Accountability}


\subsection{Reinikka, Svensson (2005): Fighting Corruption to Improve Schooling: Evidence from a Newspaper Campaign in Uganda}
\begin{transitionframe}
\begin{center}
  \begin{huge}
    \textcolor{NTHU1}{%
Reinikka, Svensson (2005):\\ Fighting Corruption to Improve Schooling: Evidence from a Newspaper Campaign in Uganda
    }
  \end{huge}
\end{center}
\end{transitionframe}


\begin{frame}
\frametitle{Q: How can monitoring improve public goods and services? } 
How do you improve provision of public goods
\begin{itemize}\setlength\itemsep{3pt}
\item Many progress in reduced-form natural and randomized experiments
\item When scaling the findings up, understanding of the entire delivery chain matters
\begin{itemize}\setlength\itemsep{3pt}
\item What institutional constraints hinder the policy decisions?
\item What incentives influence different layers of governments, agents, and beneficiaries?
\end{itemize}
\item These are required for more fruitful policy implications
\end{itemize}\medskip
This paper highlights the significance of governance structure in public services 
\begin{itemize}\setlength\itemsep{3pt}
\item Using a policy experiment from newspaper campaign in Uganda
\item The aim is to see how providing information reduces capture of public funds
\item Induce bottom-up enforcement by publishing data on educational grant usage
\item Improved primary school enrollment and performance
 \end{itemize}
\end{frame}





\begin{frame}
\frametitle{Educational grants and newspaper experiment in Uganda} 
Educational grant program aimed to support financing nonwage expenditures of schools
\begin{itemize}\setlength\itemsep{3pt}
  \item Funds disbursed by Ministry of Education to the schools on regular basis
\item On paper, it seemed to work well (central government budget data)
\item On the ground, only 20\% of the funds reached schools, many schools received nothing
\item Bulk of the grants were captured by local government officials
\end{itemize}\medskip
In response, newspaper campaign was rolled out to combat these issues in late 1990s
\begin{itemize}\setlength\itemsep{3pt}
\item Not involving enhancement of legal and financial institutions (top-down)
\item Campaigns induce citizen enforcement through increased awareness
\item Geographical variation induced by the placement of newspaper outlet
\item Knowledge of info instrumented with above geographical variation
 \end{itemize}
\end{frame}




\begin{frame}
\frametitle{Improvements since the newspaper campaign} 
 \begin{figure}
  \includegraphics[keepaspectratio, width=1\textwidth]{rs2005_t1.png}
 \end{figure}
\end{frame}





\begin{frame}
\frametitle{Data and empirical strategy} 
Sources of financial data
\begin{itemize}\setlength\itemsep{3pt}
\item Public Expenditure Tracking Survey (PETS)
\begin{itemize}\setlength\itemsep{3pt}
\item Measures how much public resources were filtered down to \textit{schools}
\item Central government data only shows how much \textit{outflow} was under educational grant
\end{itemize}
\item Also controls for time and income of the local region
\end{itemize}\medskip
Two ways of exploring effects of differential information
\begin{itemize}\setlength\itemsep{3pt}
\item IV: Share of funding reaching school instrumented with distance to newspaper outlet
\[
\Delta \text{students}_j=\alpha+\gamma_0 \Delta x_j +\gamma_1 \sigma_{2001}+\gamma_2 \Delta \hat{s}_j+e_j
\]
\item DID: Variation in newspaper circulation per district across time
\[
\text{students}_j=\alpha+\gamma  x_j +\sigma_{1995}+\sigma_{2001}+\beta_0\text{Newspaper}_j+\beta_1\text{Newspaper}_j \times D_{1995} +e_j
\]
\end{itemize}
\end{frame}


\begin{frame}
\frametitle{Positive effects of being informed on enrollment} 
 \begin{figure}
  \includegraphics[keepaspectratio, width=1\textwidth]{rs2005_t2.png}
 \end{figure}
\end{frame}


\begin{frame}
\frametitle{Evidence also consistent on DID regression} 
 \begin{figure}
  \includegraphics[keepaspectratio, width=0.8\textwidth]{rs2005_t3.png}
 \end{figure}
\end{frame}




\begin{frame}
\frametitle{Key messages and developments} 
Key message: information can improve public service provision
\begin{itemize}\setlength\itemsep{3pt}
\item Bottom-up approach meant to improve monitoring
\item Suggests that information can lead to improvements via accountability
\item Very old paper: More recent papers have found more rigorous ways to test this idea
\end{itemize}\medskip
Developments since this paper
\begin{itemize}\setlength\itemsep{3pt}
\item Randomized audits: Reduce corruption of politicians
\begin{itemize}\setlength\itemsep{3pt}
\item Revelation of audit results affect politicians reelection chances and image
\item See Ferraz, Finan (2008), Larreguy, Marshall, and Snyder Jr (2021)
\end{itemize}
\item Collective actions: Information as mediator of protests
\begin{itemize}\setlength\itemsep{3pt}
\item Usage of online surveys, distribution of mobile data etc
\item See Manacorda, Tesei (2020), Enikolopov, Marakin, Petrova (2020), Bursztyin et al (2021) 
\end{itemize}
\end{itemize}\medskip
\end{frame}




\subsection{Ferraz, Finan (2008): Exposing Corrupt Politicians: The Effects of Brazil's Publicly Released Audits on Electoral Outcomes}
\begin{transitionframe}
\begin{center}
  \begin{huge}
    \textcolor{NTHU1}{%
Ferraz, Finan (2008):\\ Exposing Corrupt Politicians: The Effects of Brazil's Publicly Released Audits on Electoral Outcomes
    }
  \end{huge}
\end{center}
\end{transitionframe}




\begin{frame}
\frametitle{Q: Do disclosing information on corruption affect accountability?} 
Information on politicians' performance is crucial for holding them accountable
\begin{itemize}\setlength\itemsep{3pt}
\item Open information induces alignment of the interests of politicians and the public
\item However, such information is not randomly assigned
\item Rather, information availability is endogenous to voter efforts and community traits
\item Furthermore, such information could be manipulated in absence of independent media
\end{itemize}\medskip
This paper studies effects of disclosing corruption practices on electoral outcomes
\begin{itemize}\setlength\itemsep{3pt}
  \item Using anticoruption program in Brazil initiated in April 2003
  \item Randomly select municipalities to be audited for use of federal funds
  \item Audit results were disseminated publicly, including to the general public
  \item Compare electoral outcomes of municipalities audited pre-election vs post-election
\end{itemize}
\end{frame}



\begin{frame}
\frametitle{Anticorruption program in Brazil} 
Rollout of anticorruption program in 2003
\begin{itemize}\setlength\itemsep{3pt}
\item Aimed at discouraging misuse of public funds and promote bottom-up monitoring
\item Municipalities subject to audits determined by lottery
\item If chosen, all info on federal funds transferred to the auditing body
\item Then, representative auditors sent to municipal governments
\begin{itemize}\setlength\itemsep{3pt}
\item Also receive complaints from community members and municipal councilmen
\end{itemize}
\item Detailed reports on irregularities submitted to the central goverment
\item Summary of findings posted on the web and disclosed to media sources
\item Most voters receives this information via local radio stations
\begin{itemize}\setlength\itemsep{3pt}
\item Newspaper subscription is low (42 copies per 1000)
\end{itemize}
\end{itemize}\medskip
\end{frame}



\begin{frame}
  \frametitle{Data and empirical strategy}
Data sources are primarily audit reports released to public 
\begin{itemize}\setlength\itemsep{3pt}
\item From 373 municipalities (municipalities ran by 1st-term mayors)
\item Total federal funds, audited amounts, itemized list of irregularities
\item Corruption: Frauds in procurements, fake receipts, embezzlement, over-invoicing
\item Combined with geographical surveys and electoral outcomes of `04 local elections
\end{itemize}\medskip
Estimation leverages variation in audit timing($A_{ms}$), controlling for \# of irregularities($C_{ms}$)
\[
E_{ms}=\alpha + \beta_0 C_{ms}+\beta_1 A_{ms}+\beta_2(A_{ms}\times C_{ms})+\gamma X_{ms}+\nu_s + e_{ms}
\]\vspace{-20pt}
\begin{itemize}\setlength\itemsep{3pt}
\item Compare municipalities audited before ($A_{ms}$=1) to those audited after
\item Control for the degree of seriousness in corruption
\item Also uses variation by media source in some specifications
\end{itemize}
\end{frame}

\begin{frame}
  \frametitle{Reduction in re-election probabilities following pre-election audits}
  \begin{figure}
    \centering
    \includegraphics[keepaspectratio, width=1\textwidth]{ff2008_t1.png}
  \end{figure}
\end{frame}

\begin{frame}
  \frametitle{Reduction in re-election probabilities following pre-election audits}
  \begin{figure}
    \centering
    \includegraphics[keepaspectratio, width=1\textwidth]{ff2008_t2.png}
  \end{figure}
\end{frame}

\begin{frame}
  \frametitle{Results evident in simple relationships}
  \begin{figure}
    \centering
    \includegraphics[keepaspectratio, width=0.7\textwidth]{ff2008_f1.png}
  \end{figure}
\end{frame}



\begin{frame}
\frametitle{Key findings and remaining questions} 
Effect of information on electoral outcomes
\begin{itemize}\setlength\itemsep{3pt}
\item Providing information on corruptions enhance accountability of politicians
\item Particularly more evident if information is spread via dominant channel (radio)
\item It also rewards well-behaving politicians
  \end{itemize}\medskip
Remaining questions
\begin{itemize}\setlength\itemsep{3pt}
\item This is a study based on one election
\item Can this effect be sustained over time?
\item What about in the environment where the mass media is not trusted?
 \begin{itemize}\setlength\itemsep{3pt}
\item Or worse: When media spread misinformation? 
\end{itemize}
\item Political accountability in new, social media environment is actively being studied
\end{itemize}
\end{frame}



\subsection{Artiles, Klein-Rueschkam, León-Cilotta (2021): Accountability, Political Capture, and Selection into Politics: Evidence from Peruvian Municipalities}
\begin{transitionframe}
\begin{center}
  \begin{huge}
    \textcolor{NTHU1}{%
Artiles, Klein-Rueschkam, León-Cilotta (2021):\\ Accountability, Political Capture, and Selection into Politics: Evidence from Peruvian Municipalities
    }
  \end{huge}
\end{center}
\end{transitionframe}


\begin{frame}
  \frametitle{Q: How does politically captured accountability mechanism function?}
Can measures meant to increase accountability work when they could be manipulated?
\begin{itemize}\setlength\itemsep{3pt}
\item In principle, many mechanisms are in place to monitor politicians action and selection
\begin{itemize}\setlength\itemsep{3pt}
\item Re-election incentives, media, recall mechanism, impeachement
\end{itemize}
\item In low capacity countries, these measures could be subject to political capture
\item Could be driven by special interest groups or manipulated by political elites
\item Can the objectives of the accountability measure be distorted?
\end{itemize}\medskip
This paper investigates how misuse of these mechanisms affects selection of politicians
\begin{itemize}\setlength\itemsep{3pt}
\item Recall elections initiated by special interest groups in Peru
\item Municipalities whose mayor barely survived recall vs ousted by small margin
\item Intrinsically motivated/high-ability politicians discouraged from running 
\begin{itemize}\setlength\itemsep{3pt}
\item Recall of mayors negatively affects net value of the position
\end{itemize}
\end{itemize}
\end{frame}

\begin{frame}
  \frametitle{Peruvian recall procedures and political landscape}
Local politics in Peru at municipal level
\begin{itemize}\setlength\itemsep{3pt}
\item Highly decentralized and responsible basic public goods (streets, security, waste)
\item Mayors on 4-year terms with possibility of re-election (not after 2015)
\item Very fragmented: Many independent parties based on candidates, not ideology
\end{itemize}\medskip
How recall elections work in Peru (Pretty common  in Peru)
\begin{itemize}\setlength\itemsep{3pt}
\item Voters can recall local elected officials at 2nd or 3rd year of the term
\item Buy recall kit (forms) $\to$ name politician $\to$ collect signature (25\% of voters)
\item Recalled if turn out is at least 50\% and if 50\% + 1 of valid votes call for recall
\item No makeup election on the spot (first councilor takes over)
\item 20,000 recall attempts in `97-`13, with 5,000 officials facing referendum
\begin{itemize}\setlength\itemsep{3pt}
\item Often losers of previous elections gathering support to oust incumbent 
\item Mainly in small cities ($\to$ policymakers in Lima were unaware) 
\item Also in places where where mayor won by small share (many win with 35-ish\%)
\end{itemize}
\end{itemize}
\end{frame}

\begin{frame}
  \frametitle{Data and empirical strategy}
Data sources
\begin{itemize}\setlength\itemsep{3pt}
\item Candidate characteristics obtained from www.infogob.com.pe
\begin{itemize}\setlength\itemsep{3pt}
\item Maximal educational level, years of schooling, previous experience, demographics
\end{itemize}
\item Electoral outcomes from national election office (ONPE)
\item Budget outcomes from Ministry of Economy and Finance (MEF)
\end{itemize}\medskip
Empirical design leverages differential returns to mayorship across individuals
\begin{itemize}\setlength\itemsep{3pt}
\item Recall lowers expected term length regardless of the candidates potential performance
\item Drive out policy-driven candidates and those with better outside options
 \begin{itemize}\setlength\itemsep{3pt}
\item High-ability individuals, those with policy experience $\downarrow$ Less-able ones $\uparrow$
\end{itemize}
\item Leverage sharp regression design based on close recall election outcomes by vote shares
\[
Y_{imt}=\alpha + \beta \text{Recalled}_{mt-1}+\gamma f(\text{voteshare}_{mt-1})+e_{imt}
\]
\end{itemize}
\end{frame}


\begin{frame}
  \frametitle{Decline in highly-educated candidates following referendum}
  \begin{figure}
    \centering
    \includegraphics[keepaspectratio, width=0.65\textwidth]{akl2021_f1.png}
  \end{figure}
\end{frame}

\begin{frame}
  \frametitle{Recalls are politically driven \& have effects regardless of performance}
  \begin{figure}
    \centering
    \includegraphics[keepaspectratio, width=1\textwidth]{akl2021_t1.png}
  \end{figure}
\end{frame}


\begin{frame}
  \frametitle{Lower quality of mayors and less revenue \& expenditures}
  \begin{figure}
    \centering
    \includegraphics[keepaspectratio, width=1\textwidth]{akl2021_t2.png}
  \end{figure}
\end{frame}


\begin{frame}
\frametitle{Key findings and remaining questions}
Key findings
\begin{itemize}\setlength\itemsep{3pt}
\item Accountability measures not only affect incumbents, but also potential candidates
\item If the procedures are captured, there could be negative effects
 \begin{itemize}\setlength\itemsep{3pt}
  \item High ability individuals discouraged from seeking positions
  \item Less representative to indigenous population
  \item More inefficient public finances 
\end{itemize}

\end{itemize}\medskip
Remaining questions
\begin{itemize}\setlength\itemsep{3pt}
\item Mechanisms designed to give citizens control over politics could be distorted
\item Particularly in cases where manipulation due to political motives are unchecked
\item How to design safeguards to prevent political capture of these procedures?
\end{itemize}
\end{frame}



\section{Incentives on bureaucrats}
\subsection{Duflo, Hanna, Ryan (2012): Incentives Work: Getting Teachers to Come to School}
\begin{transitionframe}
\begin{center}
  \begin{huge}
    \textcolor{NTHU1}{%
Duflo, Hanna, Ryan (2012):\\ Incentives Work: Getting Teachers to Come to School
    }
  \end{huge}
\end{center}
\end{transitionframe}


\begin{frame}
\frametitle{Q: How does monitoring and incentives reduce absentivism?}
Expansion of school access often not accompanied by improvements in school quality
\begin{itemize}\setlength\itemsep{3pt}
\item One major factor is teachers not showing up to work
\begin{itemize}\setlength\itemsep{3pt}
\item India: about 24 \% of teachers were absent during school hours
\item Para-teachers on temporary contracts often fill in for teachers' absence
\end{itemize}
\item Regular teachers are politically entrenched and can resist attempts to enforce rules
\item Thus, monitoring and financial incentives may work better on para-teachers
\end{itemize}\medskip
This paper tests whether monitoring and incentives reduce absenteeism of para-teachers
\begin{itemize}\setlength\itemsep{3pt}
\item Theoretically not straightforward
\begin{itemize}\setlength\itemsep{3pt}
\item More effort for high pay vs crowd out intrinsic effort or only work until target income
\end{itemize}
\item Also investigates whether schooling outcomes improve as a result
\item Randomized experiment with video device and nonlinear salaries in rural India
\item Provides models that shows how para-teachers respond to financial incentives
\end{itemize}
\end{frame}



\begin{frame}
\frametitle{Experimental setup} 
Experimental design (jointly at Udaipur organized with Seva Mandir (NGO))
\begin{itemize}\setlength\itemsep{3pt}
\item Education carried out in nonformal education centers (NFEs)
\begin{itemize}\setlength\itemsep{3pt}
\item 1 para-teacher per NFE, teaching about 20 students
\item When the para-teacher is absent, the school is closed
\end{itemize}
\item Base salary: \$ 160 PPP for at least 20 days of work per month
\item in `03, the NGO assigned 60 schools in the treatment group, 60 in control
\begin{itemize}\setlength\itemsep{3pt}
\item Teachers given camera and students instructed to take picture of teachers and students
\item Each camera came with tamper-proof settings
\item \$1.15 PPP bonus (penalty) for each additional (shortfall) day in excess of 20 days 
\end{itemize}
\end{itemize}\medskip
Data collection
\begin{itemize}\setlength\itemsep{3pt}
\item Attendance: Unannounced visits and camera data
\item School activities: \# of children, writing on blackboard, teacher talk w/ students?
\item Performance measured with 3 competency exams (pre/mid/post tests)
\end{itemize}
\end{frame}

\begin{frame}
  \frametitle{Teacher attendance markedly increases}
  \begin{figure}
    \centering
    \includegraphics[keepaspectratio, width=1\textwidth]{dhr2012_f1.png}
  \end{figure}
\end{frame}


\begin{frame}
  \frametitle{Particularly in schools with teachers below median test scores}
  \begin{figure}
    \centering
    \includegraphics[keepaspectratio, width=1\textwidth]{dhr2012_t1.png}
  \end{figure}
\end{frame}

\begin{frame}
  \frametitle{Financial incentives do not crowd out effort}
  \begin{figure}
    \centering
    \includegraphics[keepaspectratio, width=1\textwidth]{dhr2012_t2.png}
  \end{figure}
\end{frame}

\begin{frame}
  \frametitle{Students's educational performance also improved}
  \begin{figure}
    \centering
    \includegraphics[keepaspectratio, width=1\textwidth]{dhr2012_t3.png}
  \end{figure}
\end{frame}





\begin{frame}
\frametitle{Key messages and the remaining questions} 
Monitoring combined with financial incentives work
\begin{itemize}\setlength\itemsep{3pt}
\item Decrease teacher absenteeism and improve student outcomes
\item Efforts of teachers are not crowded out while in school
\item Monitoring alone not as effective as monitoring combined with financial incentives
\end{itemize}\medskip
Remaining questions
\begin{itemize}\setlength\itemsep{3pt}
\item Can this be extended to regular teachers in government schools?
\begin{itemize}\setlength\itemsep{3pt}
\item Politically powerful and likely to resist regulations
\item How do you surmount entrenched interst to institute changes?
\end{itemize}
\item Can this work in other professions and contexts?
\begin{itemize}\setlength\itemsep{3pt}
\item Similar pilot program on local nurses not effective long-run ($\because$ many exemptions)
\item Tailoring monitoring/financial incentives accordingly?
\end{itemize}
\end{itemize}\medskip
\end{frame}





\subsection{Khan, Khwaja, Olken (2019): Making Moves Matter: Experimental Evidence on Incentivizing Bureaucrats through Performance-Based Postings}
\begin{transitionframe}
\begin{center}
  \begin{huge}
    \textcolor{NTHU1}{%
Khan, Khwaja, Olken (2019):\\ Making Moves Matter: Experimental Evidence on Incentivizing Bureaucrats through Performance-Based Postings
    }
  \end{huge}
\end{center}
\end{transitionframe}



\begin{frame}
\frametitle{Q: Can incentives through posting improve organizational efficiency?} 
Governments often aim to provide incentives through nonmonetary means
\begin{itemize}\setlength\itemsep{3pt}
 \item Financial incentives are subject to regulation and fiscal constraint
 \item Spots for promotion are often limited, especially in seniority-based systems
 \item Where people are posted can function as incentives
\begin{itemize}\setlength\itemsep{3pt}
\item Bureaucrats have preferences over different work locations and move frequently
\item In practice, posting is determined by factors other than performance
\end{itemize}
\item Challenges: Diverse preferences over postings, and may not reveal truthfully
\end{itemize}\medskip
This paper asks whether assignment of bureaucrats can work as an incentive mechanism
\begin{itemize}\setlength\itemsep{3pt}
  \item Large-scale experiment in property tax department in Punjab, Pakistan
  \item Posting purely based on performance (high-performer pick first) vs status quo
  \item Tests how new assignment mechanism affects tax revenues
  \item Identify who are most likely have stronger incentives 
  \end{itemize}
\end{frame}



\begin{frame}
\frametitle{Assignment mechanism: Performance-ranked serial dictatorship}
How it works in design
\begin{itemize}\setlength\itemsep{3pt}
    \item Serial dictatorships: Individuals take turns choosing among remaining options
    \item Known to be able to uncover true preferences over fixed set of slots
    \item Early theories are silent on how the individuals should be ordered
  \end{itemize}\medskip
  How it is implemented here and what the theory predicts
\begin{itemize}\setlength\itemsep{3pt}
    \item Ordering by performance (recovery, tax assessments)
    \item Increased performance $\to$ Likelihood of getting preferrable slot $\uparrow$
    \item Incentive effects depend on own vs. others' preference, and performance distribution
    \begin{itemize}\setlength\itemsep{3pt}
    \item If indiv $i$ prefers a popular slot $j$ - more incentives (and vice versa)
    \item If $i$ is very likely or unlikely to outperform, then lower incentives
      \end{itemize}
  \end{itemize}
\end{frame}



\begin{frame}
\frametitle{Data and experimental setup}
Background information on Punjab, Pakistan (110 million population)
\begin{itemize}\setlength\itemsep{3pt}
    \item Low tax collection: low tax base, inadequate enforcement, corruption
    \item Collection unit (circle): led by 1 inspector and covers 2,000-10,000 properties
    \begin{itemize}\setlength\itemsep{3pt}
    \item They determine how property tax formula is applied and if the property is liable
    \item Size, amenities, and ease of collecting taxes vary - inspectors care about postings
  \end{itemize}
  \item 1/3 of inspectors transferred per year: Status quo mechanism is opaque
  \end{itemize}\medskip
  Data sources and the 2-year experiment
    \begin{itemize}\setlength\itemsep{3pt}
    \item Tax revenues, base, recovery rates, non-exemption rates collected at circle level
    \item Preference data from all inspectors in all circles (rank all circles within a group)
   \begin{itemize}\setlength\itemsep{3pt}
    \item Collected before the announcement of treatment assignment
  \end{itemize}
  \item Randomization at group level (10 circles) for each year by a lottery
   \begin{itemize}\setlength\itemsep{3pt}
    \item Performance based on recovery (1/2 of treatment) and tax assessments (other half)
  \end{itemize}  
\end{itemize}
\end{frame}


  
\begin{frame}
  \frametitle{Improved tax revenue following experiment}
  \begin{figure}
    \centering
    \includegraphics[keepaspectratio, width=1\textwidth]{kko2019_t1.png}
  \end{figure}
\end{frame}

\begin{frame}
  \frametitle{But the effect dies off with the repeated treatment}
  \begin{figure}
    \centering
    \includegraphics[keepaspectratio, width=1\textwidth]{kko2019_t2.png}
  \end{figure}
\end{frame}

\begin{frame}
\frametitle{Treatment effect varies with knowledge of preference and outcomes} 
\begin{table}[H]
\centering
\begin{tabular}{lcc}
\hline\hline
 & \multicolumn{2}{c}{Revenue} \\
\cmidrule(lr){2-3}
Estimated Treatment $\times$ model-predicted effort & (1) & (2) \\
\hline
Full knowledge of \textbf{P}, \textbf{y} & 0.038 & 0.042 \\
(know others' pref \& performance distribution)& (0.051) & (0.066) \\[6pt]
Random \textbf{P}, full knowledge of \textbf{y} & 0.076 & 0.176 \\
(know performance distribution, not others' pref) & (0.036) & (0.103) \\[6pt]
Assume identical \textbf{P}, full knowledge of \textbf{y} & 0.025 & 0.051 \\
(assume same pref as own, know performance distribution) & (0.027) & (0.051) \\[6pt]
Full knowledge of \textbf{P}, no knowledge of \textbf{y} & $-$0.009 & $-$0.014 \\
(know others' pref, not performance distribution) & (0.060) & (0.069) \\
\hline
Observations & 249 & 249 \\
Mean of control group & 16.268 & 16.268 \\
\hline\hline
\end{tabular}
\end{table}
\end{frame}


\begin{frame}
\frametitle{Key findings and remaining questions} 
Other key findings
\begin{itemize}\setlength\itemsep{3pt}
  \item The effect sizes are similar to the performance pay experiment 
\begin{itemize}\setlength\itemsep{3pt}
    \item Carried out in similar contexts by the same group of authors!
\end{itemize}
\item Group that increases the effort most? 
\begin{itemize}\setlength\itemsep{3pt}
    \item Case where others' performance are predictable but not preferences
    \item Uncertainty about preferences + info on how much to improve $\to$ effort to raise rank $\uparrow$ 
\end{itemize}
\end{itemize}\medskip
Remaining questions
\begin{itemize}\setlength\itemsep{3pt}
    \item Nonpecuniary incentives can work as much as pecuniary ones
    \item Both schemes may reduce flexibility over assignment
\begin{itemize}\setlength\itemsep{3pt}
    \item Sometimes, may need to send top performers where they are needed most
    \item This may be different from agents' preferences 
\end{itemize}
\item How to balance organizational goals with individual motivation?
\end{itemize}
\end{frame}


\subsection{Bertrand, Burgess, Chawla, Xu (2020): The Glittering Prizes: Career Incentives and Bureaucratic Performance}
\begin{transitionframe}
\begin{center}
  \begin{huge}
    \textcolor{NTHU1}{%
Bertrand, Burgess, Chawla, Xu (2020):\\ The Glittering Prizes: Career Incentives and Bureaucratic Performance
    }
  \end{huge}
\end{center}
\end{transitionframe}

\begin{frame}
\frametitle{Q: How does career incentives of bureaucrats affect performance?}
What motivates bureaucrats to work hard?
\begin{itemize}\setlength\itemsep{3pt}
\item We know how politicians and private sector managers affect performance
\item But what motivates bureaucrats?
\begin{itemize}\setlength\itemsep{3pt}
\item They enter through a competitive procedures (exams)
\item They are seldom fired, promoted based on seniority, and retire at fixed age
\end{itemize}
\item Meant to limit favoritism and prevent lobbying, but where does motivation kick in?
\end{itemize}\medskip
This paper investigates this question using survey and administrative data on bureaucrats
\begin{itemize}\setlength\itemsep{3pt}
\item Using examples from Indian Administrative Service (IAS)
\item What motivates IAS workers? Promotion to top positions and perks that come with it
\item Leverage natural experiment from variation in age at entry
\begin{itemize}\setlength\itemsep{3pt}
\item Those who entered old cannot reach the top $\to$ Are they less motivated?
\item Also leverage pension reform that extended retirement age by 2 years
\end{itemize}
\end{itemize}
\end{frame}


\begin{frame}
\frametitle{Who are the IAS officers?}
Elite administrative civil service of India (that shaped other Commonwealth countries)
\begin{itemize}\setlength\itemsep{3pt}
\item Competitive entry based on exams: 120 acceptances for roughly 450,000 applicants
\item Age limit: Must be 21-30 years old at entry
\item Selected officers allocated to state cadre in the subsequent year based on rules
\end{itemize}\medskip
What are their generic career progressions?
\begin{itemize}\setlength\itemsep{3pt}
\item Promotion by seniority: After 4,9, 13, 16, 25 and 30 years to reach next payscale
 \begin{itemize}\setlength\itemsep{3pt}
\item Salary jump at final payscale: 25-60\% 
\end{itemize}
\item Actual tenure required for promotion is stringent for later promotions
 \begin{itemize}\setlength\itemsep{3pt}
\item Depends on performance reviews, and availability
\end{itemize}
\item Retirement at 58 before `98 and 60 post-`98, with pensions linked to final payscale
\item Retiring at highest payscale comes with social prestige and other nonpecuniary perks
\item Those who enter older may not reach top payscale $\to$ Incentives and effectiveness $\downarrow$

\end{itemize}
\end{frame}





\begin{frame}
  \frametitle{Data and regression analysis}
Measuring bureaucratic performance
\begin{itemize}\setlength\itemsep{3pt}
\item Surveys assessing bureaucrauts from 14 states (85\% of Indian pop.) on 1-to-5 scale
\begin{itemize}\setlength\itemsep{3pt}
\item Overall rating, job effectiveness, probity, withstrand illegal pressure, pro-poor sentiment
\item From other IAS officers, state civil servants, politicians, businessmen, media, civil society
\item Subjectivity addressed with source FE, interviewer FE, and respondent FE
\end{itemize}
\item Names of IAS workers compiled at state level - covers 70\% of all officers
\item Administrative data: Gather suspension episodes, internal rankings etc
\end{itemize}\medskip
 Regression leverage variation in age of entry and timing of reforms at `98 
\begin{itemize}\setlength\itemsep{3pt}
\item  Leveraging entry age only (agent $i$, respondent $j$):\vspace{-10pt}
\[
y_{ij}= \alpha \text{age\_entry}_{i}+\beta X_i + \theta_j + e_ij
\]\vspace{-20pt}
\item Leveraging entry age and reforms\vspace{-10pt}
\[
y_{ij}= \alpha_1 \text{age\_entry}_{i}+\alpha_2 \text{age\_entry}_{i}\times\text{post98}_i +\beta X_i + \theta_j + e_ij
\]
\end{itemize}
\end{frame}

\begin{frame}
  \frametitle{Those who enter older are less likely to be perceived as effective}
  \begin{figure}
    \centering
    \includegraphics[keepaspectratio, width=0.8\textwidth]{bbcx2020_t1.png}
  \end{figure}
\end{frame}

\begin{frame}
  \frametitle{Delayed retirment increases perceived effectiveness of older entrants}
  \begin{figure}
    \centering
    \includegraphics[keepaspectratio, width=0.7\textwidth]{bbcx2020_t2.png}
  \end{figure}
\end{frame}

\begin{frame}
  \frametitle{Performance of marginal entrants (age 28,29) improves post-reform }
  \begin{figure}
    \centering
    \includegraphics[keepaspectratio, width=0.7\textwidth]{bbcx2020_f1.png}
  \end{figure}
\end{frame}


\begin{frame}
  \frametitle{Takeaways and remaining questions}
  Key findings
  \begin{itemize}\setlength\itemsep{3pt}
\item Paper also finds that older entrants working with younger colleagues perform worse
\item Older entrants with less likelihood of reaching top levels face more suspensions
\item All evidence suggest career potentials act as incentive mechanism for bureaucrats
\item Thus highlights the role of nonmonetary instruments in shaping performance
\end{itemize}\medskip
Remaining questions
  \begin{itemize}\setlength\itemsep{3pt}
\item How do civil leaders affect societal outcomes?
 \begin{itemize}\setlength\itemsep{3pt}
\item Paper is silent on economic and social outcomes of motivated bureaucrats
\item Ultimate guidance on whether bureaucratic reforms work!
\end{itemize}
  \item Roles of other nonmonetary instruments?
\begin{itemize}\setlength\itemsep{3pt}
\item Social prestige that comes with bureaucracy, for instance
\item Potentially leverage exogenous changes in the exams selecting bureacrats?
\end{itemize}
\end{itemize}
\end{frame}

\begin{frame}
\frametitle{Unanswered questions}
Citizens as a stakeholder in all of this?
\begin{itemize}\setlength\itemsep{3pt}
\item What is the best way to monitor bureaucracy (media, SNS, civil society etc)
\item What method of public organizations serves the citizens the best?
\item How do you motivate them to join public organizations (or comply to public policies)?
\end{itemize}\medskip
Ultimately, what is the best way to organize a bureaucracy?
\begin{itemize}\setlength\itemsep{3pt}
\item Typically organized in hierarchy
\item But alternatively, some are form based on tasks or events (e.g. task forces)
\item Is clear structure better (may need bureaucrats to carry difficutl tasks against will)
\item Or is providing flexibility better?
\end{itemize}\medskip
Keep an interest on Vorada's class entirely dedicated to bureaucracy next semester
\end{frame}



%%%%%%%%%%%
\end{document}
